% Part: computability
% Chapter: computability-theory
% Section: fixed-point-thm

\documentclass[../../../include/open-logic-section]{subfiles}

\begin{document}

\olfileid{cmp}{thy}{fix}
\olsection{નિશ્ચિત-બિંદુનું પ્રમેય}

અટકવાની સમસ્યા ફરી વિચારીએ. અસ્થાયી સંકેત તરીકે
$\tuple{x, y}$ માટે $\gn{\cfind{x}(y)}$ લખીએ; તેને
$\cfind{x}(y)$ના મૂલ્યનું ``નામ'' માનો. આ સંકેત વડે
અટકવાની સમસ્યા અનિર્ણેય છે તેની આપણી એક સાબિતી
બીજી રીતે કહી શકાય.

પ્રશ્ન: શું એવું સંગણનીય વિધેય $h$ છે કે દરેક $x$ અને
$y$ માટે
\[
h(\gn{\cfind{x}(y)}) =
\begin{cases}
1 & \text{જો $\cfind{x}(y) \fdefined$} \\
0 & \text{અન્યથા.}
\end{cases}
\]

જવાબ: ના. અન્યથા આંશિક વિધેય
\[
g(x) \simeq
\begin{cases}
0 & \text{જો $h(\gn{\cfind{x}(x)}) = 0$} \\
\text{અનિર્ધારિત} & \text{અન્યથા}
\end{cases}
\]
સંગણનીય હોત અને તેથી તેનો કોઈ સૂચક~$e$ હોત. પરંતુ પછી
\[
\cfind{e}(e) \simeq
\begin{cases}
0 & \text{જો $h(\gn{\cfind{e}(e)}) = 0$} \\
\text{અનિર્ધારિત} & \text{અન્યથા,}
\end{cases}
\]
એટલે $\cfind{e}(e)$ ત્યારે અને માત્ર ત્યારે વ્યાખ્યાયિત
થાય જ્યારે તે અનિર્ધારિત હોય; આ વિરોધાભાસ છે.

હવે $\cfind{e}$વાળું સમીકરણ જુઓ. તેમાં એક અર્થમાં
સ્વ-સંદર્ભ છે: $\cfind{e}(e)$નું મૂલ્ય એક ચોક્કસ રીતે
$\gn{\cfind{e}(e)}$ પર આધાર રાખે એવું આપણે ગોઠવ્યું છે.
નિશ્ચિત-બિંદુનું પ્રમેય કહે છે કે સામાન્ય રીતે પણ આ કરી
શકાય છે---માત્ર વિરોધાભાસ સાબિત કરવા માટે જ નહીં.

\olref{lem:fixed-equiv} નિશ્ચિત-બિંદુનું પ્રમેય બે સમકક્ષ
રીતે રજૂ કરે છે. તર્કની દૃષ્ટિએ બંને વિધાનો સાચાં હોવાથી
તેમની સમકક્ષતા મળે છે; પરંતુ ખરેખર કહેવાનો અર્થ એ છે કે
દરેક વિધાન બીજામાંથી સીધું મળે છે, તેથી તેમને એક જ પ્રમેયનાં
વૈકલ્પિક વિધાનો માની શકાય.

\begin{lem}
\ollabel{lem:fixed-equiv}
આગળનાં વિધાનો સમકક્ષ છે:
\begin{enumerate}
\item દરેક આંશિક સંગણનીય વિધેય $g(x,y)$ માટે એવો
  સૂચક~$e$ છે કે દરેક~$y$ માટે
\[
\cfind{e}(y) \simeq g(e,y).
\]
\item દરેક સંગણનીય વિધેય~$f(x)$ માટે એવો સૂચક~$e$
  છે કે દરેક~$y$ માટે
\[
\cfind{e}(y) \simeq \cfind{f(e)}(y).
\]
\end{enumerate}
\end{lem}

\begin{proof}
$(1) \Rightarrow (2)$: $f$ આપેલું હોય, તો $g$ને
$g(x,y) \simeq \fn{Un}(f(x),y)$ વડે વ્યાખ્યાયિત કરો.
(1) વડે એવો સૂચક~$e$ મેળવો કે દરેક~$y$ માટે
\begin{align*}
\cfind{e}(y) & = \fn{Un}(f(e),y) \\
& = \cfind{f(e)}(y).
\end{align*}

$(2) \Rightarrow (1)$: $g$ આપેલું હોય, તો $s$-$m$-$n$
પ્રમેયથી એવું $f$ મેળવો કે દરેક $x$ અને~$y$ માટે
$\cfind{f(x)}(y) \simeq g(x,y)$. (2) વડે એવો સૂચક~$e$
મેળવો કે
\begin{align*}
\cfind{e}(y) & = \cfind{f(e)}(y) \\
& = g(e,y).
\end{align*}
આથી સાબિતી પૂરી થાય છે.
\end{proof}

\begin{explain}
વિધાન (1) સાચું છે (અને તેથી (2) પણ) એ બતાવતાં પહેલાં
તે કેટલું નવાઈભર્યું છે તે વિચારો. $e$ને કમ્પ્યુટર પ્રોગ્રામ
માનો; વિધાન (1) કહે છે કે કોઈ પણ આંશિક સંગણનીય
$g(x,y)$ માટે $g_e(y) \simeq g(e,y)$ ગણતો કમ્પ્યુટર
પ્રોગ્રામ $e$ શોધી શકાય છે. બીજા શબ્દોમાં, પોતાને જ
સંદર્ભિત કરતું વિધેય ગણતો પ્રોગ્રામ શોધી શકાય છે.
\end{explain}

\begin{thm}
\olref{lem:fixed-equiv}નાં બંને વિધાનો સાચાં છે. ખાસ
કરીને, દરેક આંશિક સંગણનીય વિધેય $g(x,y)$ માટે એવો
સૂચક~$e$ છે કે દરેક~$y$ માટે
\[
\cfind{e}(y) \simeq g(e,y).
\]
\end{thm}

\begin{proof}
જરૂરી ઘટકો ઉપરની અટકવાની સમસ્યાની ચર્ચામાં ગર્ભિત છે.
$\fn{diag}(x)$ એવું સંગણનીય વિધેય હોય કે દરેક $x$ માટે
તે વિધેય $f_x(y) \simeq \cfind{x}[2](x,y)$નો સૂચક
પાછો આપે, એટલે
\[
\cfind{\fn{diag}(x)}(y) \simeq \cfind{x}[2](x,y).
\]
$\fn{diag}$ને એવી પ્રક્રિયા માનો કે જે દ્વિ-આર્ગ્યુમેન્ટવાળા
વિધેયના પ્રોગ્રામમાંથી એક-આર્ગ્યુમેન્ટવાળા વિધેયનો પ્રોગ્રામ
બનાવે છે, જ્યાં મૂળ પ્રોગ્રામ પોતે પ્રથમ આર્ગ્યુમેન્ટ તરીકે
નિશ્ચિત કરાયો છે. $\fn{diag}$ને ઔપચારિક રીતે આ રીતે
વ્યાખ્યાયિત કરી શકાય: પહેલાં $s$ને આમ વ્યાખ્યાયિત કરો:
\[
s(x,y) \simeq \fn{Un}^2(x,x,y),
\]
જ્યાં $\fn{Un}^2$ દ્વિ-આર્ગ્યુમેન્ટવાળાં આંશિક સંગણનીય
વિધેયો માટે સર્વવ્યાપી ત્રિ-આર્ગ્યુમેન્ટવાળું વિધેય છે.
પછી $s$-$m$-$n$ પ્રમેયથી એવું આદિમ પુનરાવર્તી વિધેય
$\fn{diag}$ મળે છે કે
\[
\cfind{\fn{diag}(x)}(y) \simeq s(x,y).
\]
\sourcecorrection{OLCT-011}{સ્થિર અંગ્રેજી મૂળ આ
બાંધકામમાં બે જગ્યાએ $\cfind{x}(x,y)$ અને નીચેની
ગણતરીમાં $\cfind{\gn{l}}(\gn{l},y)$ લખે છે. બંને
દ્વિ-આર્ગ્યુમેન્ટવાળાં વિધેયો છે; અહીં તે મુજબ
$\cfind{x}[2]$ અને $\cfind{\gn{l}}[2]$ રાખ્યાં છે.}

હવે વિધેય $l$ને આમ વ્યાખ્યાયિત કરો:
\[
l(x,y) \simeq g(\fn{diag}(x),y).
\]
અને $\gn{l}$ એ $l$નો સૂચક હોય. છેલ્લે
$e = \fn{diag}(\gn{l})$ રાખો. ત્યારે દરેક $y$ માટે
\begin{align*}
\cfind{e}(y) & \simeq \cfind{\fn{diag}(\gn{l})}(y) \\
& \simeq \cfind{\gn{l}}[2](\gn{l}, y) \\
& \simeq l(\gn{l}, y) \\
& \simeq g(\fn{diag}(\gn{l}),y) \\
& \simeq g(e, y),
\end{align*}
જે જોઈતું હતું.
\end{proof}

\begin{explain}
અહીં શું બને છે? માનો કે પોતાને જ છાપતો કમ્પ્યુટર પ્રોગ્રામ
લખવાનું કામ તમને મળ્યું છે. વધુમાં માનો કે તમે એવી
પ્રોગ્રામિંગ ભાષામાં કામ કરો છો જેમાં શબ્દમાળાઓનાં વિપુલ
અને વિચિત્ર વિધેયોનું ગ્રંથાલય છે. ખાસ કરીને, તેમાં
$\fn{diag}$ નામનું વિધેય હોય જે આ રીતે કામ કરે:
ઇનપુટ શબ્દમાળા~$s$ આપવામાં આવે, તો $\fn{diag}$ એ~$s$માં
આવતાં `x' પ્રતીકનાં દરેક સ્થાનને શોધીને તેના સ્થાને
મૂળ શબ્દમાળાનું અવતરણ-ચિહ્નવાળું
રૂપ મૂકે છે. દાખલા તરીકે, આ શબ્દમાળા આપો:
\begin{quote}
\begin{verbatim}
hello x world
\end{verbatim}
\end{quote}
તો તે પાછું આપે:
\begin{quote}
\begin{verbatim}
hello 'hello x world' world
\end{verbatim}
\end{quote}
આ સ્થિતિમાં ઇચ્છિત પ્રોગ્રામ લખવો સહેલો છે; તપાસો કે
\begin{quote}
\begin{verbatim}
print(diag('print(diag(x))'))
\end{verbatim}
\end{quote}
કામ કરે છે. C++ અને Java જેવી સામાન્ય ભાષાઓમાં પણ
આ જ વિચાર (વધુ જટિલ અમલ સાથે) કામ કરે છે.

નિશ્ચિત-બિંદુના પ્રમેયની સાબિતીથી હવે માત્ર થોડાં પગલાં
દૂર છીએ. માનો કે છાપવાના વિધેયનું એક રૂપ
$\fn{print}(x,y)$ શબ્દમાળા $x$ અને સંખ્યાત્મક આર્ગ્યુમેન્ટ
$y$ લે છે અને શબ્દમાળા $x$ને $y$ વાર છાપે છે. ત્યારે આ
``પ્રોગ્રામ''
\begin{quote}
\begin{verbatim}
getinput(y);
print(diag('getinput(y); print(diag(x), y)'), y)
\end{verbatim}
\end{quote}
ઇનપુટ $y$ પર પોતાને $y$ વાર છાપે છે.
$\fn{getinput}$---$\fn{print}$---$\fn{diag}$ના માળખાને
કોઈ પણ વિધેય $g(x,y)$ વડે બદલીને
\begin{quote}
\begin{verbatim}
g(diag('g(diag(x), y)'), y)
\end{verbatim}
\end{quote}
મળે છે; આ પ્રોગ્રામ ઇનપુટ $y$ પર, પોતે અને $y$
આર્ગ્યુમેન્ટો લઈને $g$ ચલાવે છે. ``અવતરણમાં મૂકવા''ને
``સૂચક વાપરવા'' તરીકે સમજીએ, તો ઉપરની સાબિતી મળે છે.

હમણાં માટે, આ સાબિતી ઔપચારિક યુક્તિ અથવા જાદુ જેવી
લાગે તો ચાલે. પરંતુ ઉપરની દલીલની વિગતો તમે ફરીથી
ઘડી શકો એવું હોવું જોઈએ. અપૂર્ણતાનાં પ્રમેયો (અને
સંબંધિત ``નિશ્ચિત-બિંદુનું પ્રમેય'') સાબિત કરીએ ત્યારે
આ રીત કેમ કામ કરે છે તે સમજવાની બીજી રીતો પણ જોઈશું.
\end{explain}

\begin{tagblock}{lambda}
\begin{digress}
આ જ વિચારથી ``નિશ્ચિત-બિંદુ'' સંયોજક મેળવી શકાય છે.
માનો કે તમારી પાસે લૅમ્બ્ડા પદ $g$ છે અને એવું બીજું
પદ $k$ જોઈએ છે કે $k$, $gk$ સાથે $\beta$-સમકક્ષ હોય.
આપણી સંકેત-પદ્ધતિથી પદો વ્યાખ્યાયિત કરો:
\[
\fn{diag}(x) = xx
\]
અને
\[
l(x) = g(\fn{diag}(x))
\]
બીજા શબ્દોમાં, $l$ એ પદ $\lambd[x][g(xx)]$ છે. $k$
એ પદ $ll$ હોય. ત્યારે
\begin{align*}
k & = (\lambd[x][g(xx)])(\lambd[x][g(xx)]) \\
& \red  g((\lambd[x][g(xx)])(\lambd[x][g(xx)])) \\
& = gk.
\end{align*}
જો
\[
Y = \lambd[g][((\lambd[x][g(xx)])(\lambd[x][g(xx)]))]
\]
રાખીએ, તો $Yg$ અને $g(Yg)$ બંને એક સામાન્ય પદ સુધી
ઘટે છે; એટલે $Yg \equiv_\beta
g(Yg)$. તેને ``કરીનો સંયોજક'' કહે છે. તેના બદલે જો
\[
Y = (\lambd[xg][g(xxg)])(\lambd[xg][g(xxg)])
\]
રાખીએ, તો ખરેખર $Yg$, $g(Yg)$ સુધી ઘટે છે, જે વધુ
મજબૂત વિધાન છે. $Y$ના આ બીજા રૂપને ``ટ્યુરિંગનો
સંયોજક'' કહે છે.
\end{digress}
\end{tagblock}

\end{document}
